%------------------------------------------------------------------------------%

\usepackage{calculo_varias_variaveis-1}

\title{Polinômios em Números Complexos}

%------------------------------------------------------------------------------%
\begin{document}

\addtitle

%------------------------------------------------------------------------------%
\section{Polinômios}

%------------------------------------------------------------------------------%
\begin{frame}{Polinômios}
  Um \emph{polinômio} de grau $n\in\N$ de uma variável complexa $z$ \pause é
  \[
    P_n(z)
    = a_n z^n
    + a_{n-1} z^{n-1}
    + \cdots
    + a_2 z^2
    + a_1 z
    + a_0
  \]
  com \(z\), \(a_k\in\mathbb{C}\)

  \pause
  \vspace{\baselineskip}
  Um polinômio de grau $n$ possui $n$ raízes complexas

  \pause
  \vspace{\baselineskip}
  Se os coeficientes \(a_k\) forem reais, as raízes serão reais ou pares conjugados
\end{frame}

%------------------------------------------------------------------------------%
\begin{frame}{Raízes da Equação Quadrática}
  Raízes do polinômio
  \[
    ax^2 + bx + c = 0
    \qquad\qquad
    a\neq 0
    \qquad
    a, b, c \in \R
  \]

  \pause
  Fórmula de Bhaskara
  \[
    x = \dfrac{\displaystyle -b \pm \sqrt{\Delta}}{2a}
    \qquad\text{ onde }\quad
    \Delta = b^2 - 4 ac
  \]

  \pause
  Se \(\;\Delta<0\;\) temos
  \(\quad \sqrt{\Delta} = \sqrt{\abs{\Delta}}\;i \)
\end{frame}

%------------------------------------------------------------------------------%
%------------------------------------------------------------------------------%
\begin{question}
  Encontre as raízes do polinômio
  \[
    p = 2z^2 -4z + 4
  \]
\end{question}

%------------------------------------------------------------------------------%
\begin{resolution}{Solução}
\begin{columns}
\column{0.5\linewidth}
  \begin{align*}
	  \onslide<+->{\Delta        & = b^2 - 4 ac                   \\}
	  \onslide<+->{              & = (-4)^2 - 4 \times 2 \times 4 \\}
	  \onslide<+->{              & = 16 - 32                      \\}
	  \onslide<+->{              & = -16                          \\[5mm]}
	  \onslide<+->{\sqrt{\Delta} & = \sqrt{-16}                   \\}
	  \onslide<+->{              & = \sqrt{16}\,i                 \\}
	  \onslide<+->{              & = 4i}
  \end{align*}
\column{0.5\linewidth}
  \begin{align*}
	  \onslide<+->{z & = \frac{-b \pm \sqrt{\Delta}}{2a} \\}
	  \onslide<+->{  & = \frac{-(-4) \pm 4i}{2\times 2}  \\}
	  \onslide<+->{  & = \frac{4 \pm 4i}{4}              \\}
	  \onslide<+->{  & = 1 \pm i}
  \end{align*}
\end{columns}
\end{resolution}

%------------------------------------------------------------------------------%

%------------------------------------------------------------------------------%
\begin{question}
  Encontre as raízes do polinômio \(p(x) = x^3 -6x^2 + 10x\)
\end{question}

%------------------------------------------------------------------------------%
\begin{resolution}{Solução}
  Queremos resolver a equação
  \[
    x^3 -6x^2 + 10x = 0
  \]
  que podemos escrever como
  \[
    x\left(x^2 -6x + 10\right) = 0
  \]
  Uma das raízes é \(x=0\), usamos Bhaskara para encontrar as demais
\end{resolution}

%------------------------------------------------------------------------------%
\begin{resolution}{Solução}
  \[
    \Delta
    = b^2 - 4ac
    = (-6)^2 - 4 \times 1 \times 10
    = 36 - 40
    = -4
  \]
  \[
    \sqrt{\Delta}
    = \sqrt{-4}
    = i\sqrt{4}
    = 2i
  \]
  \[
    x
    = \frac{-b\pm\sqrt{\Delta}}{2a}
    = \frac{-(-6)\pm 2i}{2}
    = 3\pm i
  \]
\end{resolution}

%------------------------------------------------------------------------------%
\begin{resolution}{Solução}
  Portanto as raízes de $p$ são
  \[
    x_1 = 0 \qquad\qquad
    x_2 = 3-i \qquad\qquad
    x_3 = 3+i
  \]
\end{resolution}

%------------------------------------------------------------------------------%


%------------------------------------------------------------------------------%
\section{Lista Mínima}

%------------------------------------------------------------------------------%
\begin{frame}{Lista Mínima}
%%  Cálculo Vol. 2 do Thomas 12\textsuperscript{a} ed. -- Seção 14.8
%%  \begin{enumerate}
%%    \item Estudar o texto da seção
%%    \item Resolver os exercícios: 3, 5, 7, 12, 16, 18, 24 e 30
%%  \end{enumerate}

  \vfill
  Atenção: A prova é baseada no livro, não nas apresentações
\end{frame}

%------------------------------------------------------------------------------%
\end{document}
%------------------------------------------------------------------------------%
